Authors of interactive math texts sometimes employ interactivity to explain mathematical formulas, for instance to convey, step through, and visualize computations. Making formulas interactive in these ways, however, is unnecessarily effort-intensive with standard tools. To lower the threshold - "the delta" - and raise the ceiling of formula explanations, we design and implement a domain-specific language for interactive formula explanation called Delta.  It gives authors new facilities to imbue formulas with computability, explain them with step-by-step walkthroughs, and link manipulation across formulas, texts, and visuals. It is designed to integrate with LaTeX math in web-based writing contexts, with a runtime that handles necessary state management. The DSL's constructs are grounded in a review of formulas in public, polished interactive math texts. In two studies, we show that the DSL makes authoring of simple interactive formulas faster and easier; we also thoroughly document the fit of the language's numerous constructs to the task. 